\documentclass[12pt,a4paper,article]{memoir}
%
\usepackage{amsmath,amssymb,amsthm}
\usepackage{hyperref}
\usepackage{cleveref}
\usepackage[dvipsnames]{xcolor}
\usepackage{mathpazo}
\usepackage[T1]{fontenc}
\linespread{1.05}
\date{}
\title{Oligo-poly\\
\vspace{10pt}
\large
A library for working with polynomial rings under oligomorphic group actions.
$ $\\}
%
\author{Arka Ghosh}
%
\newcommand{\struct}{\mathcal{X}}
\newcommand{\dloQ}{\mathcal{Q}}
\newcommand{\ordZ}{\mathcal{Z}}

\newcommand{\N}{\mathbb{N}}
\newcommand{\Q}{\mathbb{Q}}
\newcommand{\F}{\mathbb{F}}

\newcommand{\aut}[1]{\mathsf{Aut}\left(#1\right)}
\newcommand{\poly}[2]{#1[#2]}
\newcommand{\setof}[2]{\left\{#1\ \mid\ #2 \right\}}
\newcommand{\idGen}[1]{\langle#1\rangle}
\newcommand{\eqIdGen}[1]{\langle#1\rangle_{\textnormal{eqv}}}

\newcommand{\arka}[1]{}
%{\textcolor{red}{Arka:#1}}
\newcommand{\emphdef}[1]{\emph{\textcolor{BrickRed}{#1}}}
\newcommand{\code}[1]{\texttt{\textcolor{magenta}{#1}}}

\newtheorem{theorem}{Theorem}
\newtheorem{lemma}[theorem]{Lemma}

\theoremstyle{definition}
\newtheorem{example}[theorem]{Example}
\newtheorem{notation}[theorem]{Notation}
\newtheorem{todo}[theorem]{Todo}
\newtheorem{remark}[theorem]{Remark}
%
\begin{document}
\maketitle
\raggedbottom
%
\tableofcontents
%
\pagebreak
\noindent
The reader is suggested to look at the most recent version of this document which can be found in the url below.
\[
\text{
\url{https://github.com/arkagh0sh/oligo-poly}
}
\]
The purpose of this library is to implement (in Rust) the equivariant Buchberger algorithm given in \cite{GL25}.
This document is mostly self-complete in terms of definitions.
%
\paragraph*{Structure of this document}
%
We start with explaining the theoretical background in \Cref{sec:theory}.
%\Cref{sec:how} gives basic instructions for using this library.
\Cref{sec:impl} describes the important aspects of our implementation and gives basic instructions for usage.
%
\chapter{Theoretical background}\label{sec:theory}
%

\section{Relational structures}
%
A \emphdef{relational structure} $\struct = (X,R_1,\dots,R_n)$ consists of a
set $X$ (called the \emphdef{domain} of the relational structure) and relations $R_i$ defined on the set $X$.
%
\begin{example}
Let $Q$ be the set of rational numbers.
We can consider the structure $\dloQ = (Q,<)$ where is $<$ the usual ordering.
\end{example}
%
\begin{example}
A graph $\mathcal{G} = (V,E)$ is another example of a relational structure,
with the set $V$ of vertices being the domain and the edge relation $E$ being the only relation.
\end{example}

By abuse of notation we sometime use $\struct$ to denote the set underlying $X$.
For instance,
we write $x \in\struct^n$ to mean that $x$ is an element of $X^n$.

\section{Automorphism groups}

An \emphdef{automorphism} of the relational structure $\struct$ is a bijection from $\struct$ to itself that preserves and reflects the structure.
%
\begin{example}
An automorphism of $\dloQ$ is an order preserving bijection from $\mathcal{Q}$ to itself.
\end{example}
%
\begin{example}
If $\struct$ is a graph,
then an automorphism of $\struct$ is an isomorphism of the graph with itself.
\end{example}
%
The set of automorphisms of a relational structure $\struct$ is denoted as \emphdef{$\mathbf{Aut}(\struct)$}.
The group $\aut{\struct}$ acts on $\struct^n$ coordinatewise:
$\pi(x_1,\dots,x_n) = (\pi(x_1),\dots,\pi(x_n))$.
The \emphdef{orbit} of an element $x$ of a set on which $\aut{\struct}$ acts is the set $\setof{\pi(x)}{\pi\in\aut{\struct}}$.
%
\begin{example}
In $\dloQ$, the orbit of $(1,2)\in\dloQ^2$ is $\setof{(a,b)\in\dloQ^2}{a<b}$.
\end{example}
%
\section{Oligomorphic structures}
%
A relational structure $\struct$ is called \emphdef{oligomorphic} if for every $n\in\N$,
the set $\struct^n$ splits into finitely many orbits under the action of $\aut{\struct}$.
%
\begin{example}
The structure $\dloQ$ is oligomorphic.
For illustration, for $n = 2$, $\dloQ^2$ splits into three orbits
\[
\setof{(a,b)\in\dloQ^2}{a < b},
\setof{(a,b)\in\dloQ^2}{a = b}
\text{ and }
\setof{(a,b)\in\dloQ^2}{a > b} \ .
\]
\end{example}
%
\begin{example}
The structure $\ordZ = (Z,<)$ of integers with order is not oligomorphic.
For instance,
the set $\ordZ^2$ splits into infinite many orbits of the form
\[
\setof{(a,b)}{a - b = d}
\]
for $d\in Z$.
\end{example}
%
\section{Polynomial ideals}
%
We use $\poly{\F}{X}$ to denote the polynomial rings with coefficients from the field $\F$ and variables from the set $X$.
A subset $J$ of polynomials is called an \emphdef{ideal} if
\begin{enumerate}
\item the zero polynomial $0$ is inside $J$,
\item the sum $f + g$ of any two polynomials $f$ and $g$ inside $J$ is also inside $J$, and
\item the product $h\cdot f$ of any polynomial $h$ in the ring $\poly{\F}{X}$ with any polynomial $f\in J$ is also in $J$.
\end{enumerate}
%
\begin{notation}
For a function $\alpha : X\to \F$ and $f\in\poly{\F}{X}$ let $f(\alpha)$ denote the instantiation of $f$ at $\alpha$.
For instance, if $f = x^2 - y$, $\alpha(x) = 1$ and $\alpha(y) = 1.5$,
then $f(\alpha) = -0.5$.
\end{notation}
%
\begin{example}
For any subset $S\subseteq X\to \F$,
the set
\[
\setof{f\in\poly{\F}{X}}{f(\alpha) = 0\text{ for all }\alpha\in S}
\]
is an ideal inside $\poly{\F}{X}$.
\end{example}
%
\section{Equivariant ideals}
%
For a relational structure $\struct$,
we use $\mathbb{Q}[\struct]$ to denote the polynomial ring whose variables are elements of $\struct$.

Instead of using $\struct$ as the set of variables some articles use variables $a_{x}$ indexed by elements $x\in\struct$.
But using $\struct$ itself makes the notation simpler.

The group $\mathbf{Aut}(\struct)$ of automorphisms of $\struct$ acts on $\mathbb{Q}[\struct]$ in a variable-wise manner.
For instance, for $\pi\in\mathbf{Aut}(\struct)$ and $x,y\in\struct$
$$\pi(x^2 - 2xy^2 + 3) =
\pi(x)^2 - 2\pi(x)\pi(y)^2 + 3\ .$$

An ideal in $\poly{\Q}{\struct}$ is called \emphdef{equivariant},
if it is invariant under the action of $\aut{\struct}$.
%
\begin{example}
The set $J_0$ of polynomials in $\poly{\Q}{\struct}$ whose coefficients add up to zero is an equivariant ideal.

Equivalently, $J_0 = \setof{f\in\poly{\Q}{\struct}}{f(\mathbf{1}) = 0}$ where $\mathbf{1}$ is the function $\struct\to\Q$ defined as $\mathbf{1}(a) = 1$ for all $a\in\struct$.
\end{example}
%
For a subset $G\subseteq\poly{\F}{\struct}$,
the \emphdef{ideal $\idGen{G}$ generated by $G$} is the smallest ideal containing $G$.
For instance, $\idGen{\setof{a - 1}{a\in\dloQ}} = J_0$

Similarly, for a subset $G\subseteq\poly{\F}{\struct}$,
the \emphdef{equivariant ideal $\eqIdGen{G}$ generated by $G$} is the smallest equivariant ideal containing $G$.
For instance $\eqIdGen{\{a - 1\}} = J_0$ for any $a\in\dloQ$.
%
An (equivariant) ideal is called \emphdef{finitely generated} if it is generated by a finite set.

Note that equivariant makes finite generation easier.
For instance, one can prove that $J_0$ is not finite generated as an ideal,
but since $J_0 = \eqIdGen{\{a - 1\}}$,
it is finitely generated as an equivariant ideal.
%
\section{Equivariant Gröbner bases}
%
We say $\struct$ has the \emphdef{Hilbert basis property} if for every field $\F$, every equivariant ideal $J$ in $\poly{\F}{\struct}$ is finitely generated as an equivariant ideal, that is, $J = \eqIdGen{G}$ for a finite set $G$.

Theorems 11 and 12, and Corollary 13 of \cite{GL24} characterised the Hilbert's basis property of relational structures.

\arka{TODO:expand starting from this part}

We call relational $\struct$ to be \emphdef{ordered} if one of the relations of  $\struct$ induces a linear order on $\struct$.
If we restrict our attention to ordered relational structure,
then the results of \cite{GL24} collapse to the following theorem.
%
\begin{theorem}
A ordered relational structure satisfies the Hilbert's basis property if and only if finite $\N$-labelled substructures of $\struct$ are well-quasi-ordered under embeddability.
\end{theorem}
%
The article \cite{GL25} builds on their results by defining Gr\"{o}bner bases of equivariant ideals and observes that if a relational structure $\struct$ satisfies the Hilbert's basis property,
then every equivariant ideals in $\poly{\F}{\struct}$ have finite Gr\"{o}bner bases.
The central result result of \cite{GL25} is to give a modified version of the equivariant extension of the Buchberger's algorithm that takes as input a finite set $H$ of polynomials in $\poly{\Q}{\struct}$ and outputs a Gr\"{o}bner basis of the equivariant ideal generated by $H$.
They show that this algorithm terminates when $\struct$ is a ordered relational structure satisfying the Hilbert's basis property.
%
%%
%This chapter is divided into two sections.
%The first section describes how to define and work with relational structures.
%%
%\section{Relational structures}\label{sec:how str}
%
%\subsection{Pre-defined relational structures}\label{sec:pre def str}
%
%\subsubsection{Dense linear order without endpoints}\label{sec:how dlo}
%
%\subsection{Defining new relational structures}\label{sec:how new str}
%
%
%\section{Polynomial rings}
%
%\subsection{Ideals}
%
%\subsection{Buchberger's algorithm}
%
%\subsection{Equivariant Buchberger's algorithm}
%
\chapter{Implementation details}\label{sec:impl}
%
The implementation is divided into three parts.
The first part gives a \emph{framework} to write equivariant algorithms.
The second part is a simple implementation of polynomial rings $\poly{\F}{X}$ over a field $\F$,
where the variables can be of any ordered data type.
The third part combines these two and implements the equivariant Buchberger's algorithm.
\Cref{sec:trees} describes the atoms of the dense meet-tree.

The third part is described in \cref{sec:eq buchberger}.
%
\section{Building and testing}
%
The polynomial rings and the Gr\"{o}bner bases are built and tested with
\begin{verbatim}
cargo build
cargo test
\end{verbatim}
The framework for equivariant algorithms is only built with the Cargo feature \code{equivariant},
because it depends on the crate \code{rug}, which compiles the GMP library:
\begin{verbatim}
cargo test --features equivariant
\end{verbatim}
The example \code{groebner\_stdin} reads polynomials from the standard input
and prints the reduced Gr\"{o}bner basis of the ideal that they generate.
For instance,
\begin{verbatim}
echo "x*y; x + y*z" | cargo run --example groebner_stdin -- f2
\end{verbatim}
prints the reduced Gr\"{o}bner basis $\{x^2, xy, x + yz\}$ over the field with two elements.
The script \code{scripts/compare\_with\_sympy.py} compares the output of this example with SymPy on random inputs.
%
\section{Framework for equivariant algorithms}
%
We are heavily influenced by \code{LOIS} and \code{N$\lambda$}.
The main difference is that both \code{LOIS} and \code{N$\lambda$} guarantees that any program written in them is equivariant by construction.
We only guarantee that any program written following our framework is equivariant.
The benefit of doing this is that the codebase is significantly smaller and simpler.
\begin{todo}
Implement a checker that checks whether a piece of code follows the framework.
\end{todo}
%
Like \code{LOIS} and \code{N$\lambda$} we categories objects into two categories,
\emphdef{elements of orbits} (encoded by the \code{trait} \code{ElementsWithAtoms}) and \emphdef{orbits} (encoded by the \code{struct} \code{Orbit}).
\footnote{Obviously, orbits themselves can be elements of orbits.
However, since we only deal with equivariant orbits as of now,
orbits of orbits are singleton.
\arka{This is schizophrenic}}
%
\subsection{Supports and equivariance}
%
In this chapter, the relational structure $\struct$ is ordered,
that is, one of its relations is a linear order on $\struct$.
We call the elements of $\struct$ \emphdef{atoms}.
The group $\aut{\struct}$ acts on finite sets of atoms, on pairs, on lists and on polynomials
by applying the automorphism to every atom that occurs.

The \emphdef{support} $\mathsf{supp}(b)$ of an element $b$ is the set of atoms that occur in $b$.
A finite set $S$ of atoms \emphdef{supports} $b$ if $\pi(b) = b$ for every automorphism $\pi$ that fixes every element of $S$.
For all the elements in the implementation,
the support $\mathsf{supp}(b)$ supports $b$,
and $\mathsf{supp}(\pi(b)) = \pi(\mathsf{supp}(b))$ for every automorphism $\pi$.
%
\begin{example}
Let $a < b < c$ be atoms of $\dloQ$.
The support of the polynomial $a^2c + b$ is $\{a,b,c\}$,
and the support of the pair $(a,\{b,c\})$ is also $\{a,b,c\}$.
\end{example}
%
A function $f$ is \emphdef{equivariant} if $f(\pi(b)) = \pi(f(b))$ for every automorphism $\pi$ and every $b$ in the domain of $f$.
A predicate $p$, that is, a function with values true and false, is \emphdef{equivariant} if $p(\pi(b)) = p(b)$ for every automorphism $\pi$ and every $b$ in the domain of $p$.
%
\begin{example}
On pairs of atoms of $\dloQ$,
the function $(a,b)\mapsto a$ and the predicate $(a,b)\mapsto (a = b)$ are equivariant.
The constant function $(a,b)\mapsto 1$, where $1$ is the rational number one, is not equivariant,
because some automorphism $\pi$ satisfies $\pi(1) \neq 1$.
\end{example}
%
Let $S$ and $T$ be finite sets of atoms such that some automorphism $\pi$ maps $S$ onto $T$.
Since $\pi$ preserves the order,
$\pi$ maps the $i$-th smallest element of $S$ to the $i$-th smallest element of $T$, for every $i$.
So all automorphisms that map $S$ onto $T$ agree on $S$.
Their common restriction to $S$ is the \emphdef{partial automorphism} from $S$ onto $T$
(\code{PartialAut} in the code).
%
\begin{example}
In $\dloQ$, the partial automorphism from $\{1,2\}$ onto $\{5,9\}$ maps $1$ to $5$ and $2$ to $9$.
\end{example}
%
\begin{remark}\label{rem:paut}
The image $\pi(b)$ of an element $b$ only depends on the restriction of $\pi$ to $\mathsf{supp}(b)$.
Indeed, if two automorphisms $\pi$ and $\pi'$ agree on $\mathsf{supp}(b)$,
then $\pi'^{-1}\pi$ fixes every element of $\mathsf{supp}(b)$.
Since $\mathsf{supp}(b)$ supports $b$, we get $\pi'^{-1}\pi(b) = b$, that is, $\pi(b) = \pi'(b)$.
So $\pi(b)$ can be computed from the partial automorphism from $\mathsf{supp}(b)$ onto $\pi(\mathsf{supp}(b))$.
\end{remark}
%
\subsection{\code{trait AtomsWithOrd}}
%
This gives the framework for working with \emph{computationally oligomorphic} ordered relational structures.
To make a \code{struct A} satisfy this \code{trait},
one has to implement two basic operations.
\begin{enumerate}
\item \code{in\_same\_orbit(S, T)}: This function checks whether some automorphism maps the finite set \code{S : BTreeSet<A>} onto the finite set \code{T : BTreeSet<A>}.
\item \code{orbits\_in\_product(S, T)}: This function returns one representative $(S', T')$ for each orbit of pairs in the product of the orbit of \code{S} and the orbit of \code{T}.
\end{enumerate}
%
\begin{example}
In $\dloQ$, two finite sets are in the same orbit if and only if they have the same size,
because every order-preserving bijection between finite subsets of $\dloQ$ extends to an automorphism.
For $S = T = \{1\}$, the function \code{orbits\_in\_product} returns the three pairs
$(\{1\},\{1\})$, $(\{1\},\{2\})$ and $(\{2\},\{1\})$.
They represent the three orbits of pairs $(\{a\},\{b\})$ with $a = b$, $a < b$ and $a > b$.
\end{example}
%
\subsubsection{Defining relational structures}
%
One must be careful while defining a new \code{struct} as a relational structure.
Our recommended way of defining a relational structure is to use a wrapper class so the actual value can be accessed by the functions that are defined in the struct.

For illustration we define \code{DLO} as
\color{magenta}
\begin{verbatim}
#[derive(Clone, PartialEq, Eq, PartialOrd, Ord, Hash)]
pub struct DLO {
    val : Rational
}
\end{verbatim}
\color{black}
%
Observe that the field \code{val} is private.
Therefore one cannot write codes that uses operations other than the order and equality on \code{DLO}.
For instance, this disallows code snippets such as
\color{magenta}
\begin{verbatim}
let r : Rational = (x.val + y.val)/2;
let z : DLO = DLO::new(r);
\end{verbatim}
\color{black}
%
\arka{one can still write a piece of code that runs on a while loop and computes an element greater than a given element}
%
\subsection{\code{trait ElementsWithAtoms<A>}}
%
This is the \code{trait} for elements built from the atoms \code{A}.
To implement this \code{trait} for a \code{struct B} one has to implement the following functions.
\begin{enumerate}
\item \code{support()}: This computes the support of an element, as a \code{BTreeSet<A>}.
\item \code{apply\_paut(paut)}: This applies a partial automorphism whose domain contains the support of the element.
\end{enumerate}
The \code{trait} then provides the following functions.
\begin{enumerate}
\item \code{tuples\_in\_same\_orbit(first, second)}: This checks whether some automorphism maps the list \code{first} to the list \code{second}.
By \Cref{rem:paut}, it suffices to try the partial automorphism from the support of \code{first} onto the support of \code{second}.
\item \code{project\_to\_support(rep, S)}: This outputs all elements in the orbit of \code{rep} whose support is contained in \code{S : BTreeSet<A>}.
By \Cref{rem:paut}, these are the images of \code{rep} under the partial automorphisms from the support of \code{rep} onto subsets of \code{S}.
\end{enumerate}
The \code{trait} is implemented for atoms (\code{DLO}), finite sets of atoms (\code{BTreeSet<A>}),
pairs \code{(B, C)} and lists \code{Vec<B>}.
%
\subsection{\code{struct Orbit<A,B>}}
%
This is the structure of orbits of elements of another \code{struct B} under the action of the automorphism group of another \code{struct A}.
An object of the \code{struct} \code{Orbit} is given by its representative \code{rep},
which must satisfy the \code{trait} \code{ElementsWithAtoms}.
The \code{rep} field is hidden.
This ensures that orbits can only be accessed in specific ways (explained below) that preserve equivariance.

Orbits are considered as lists indexed by elements of
\code{BTreeSet<A : AtomsWithOrd>}.
This is consistent with the theoretical result that every orbit of set with atoms $\struct$ is isomorphic to an equivariant subset of $\binom{\struct}{n}$ for some $n\in\N$.
Orbits can be accessed in the following ways:
\begin{enumerate}
\item Given two orbits $B$ and $C$, one can check if they are equal (\code{==}).
\item Given an orbit $B$ and \code{S : BTreeSet<A>}, one can compute the elements of $B$ whose support is contained in \code{S} (\code{supported\_by}).
\item Given two orbits $B$ and $C$, one can compute the list of orbits in the product $B\times C$ (\code{product}).
\item Given an equivariant function $f$ with domain containing $B$,
one can compute $f(B)$ (\code{map}).
\item Given an equivariant predicate $f$ with domain containing $B$,
one can filter $B$ using $f$ (\code{filter}).
Observe that since $f$ is equivariant, this means the result of the filtering is either $B$ or $\emptyset$.
\end{enumerate}
The reason behind doing the above is the following.
First, we ensure that any function with input and output types satisfying the trait \code{ElementsWithAtoms} is equivariant.
Second, we ensure that orbits can only be accessed in specific ways that preserves equivariance.
As a consequence, any function that takes orbits as input must be equivariant.

The function \code{product} relies on the following lemma.
It computes, for each representative pair $(S',T')$ that \code{orbits\_in\_product} returns for the supports of the representatives $b$ and $c$,
the pair of images of $b$ and $c$ under the partial automorphisms onto $S'$ and $T'$.
%
\begin{lemma}\label{lem:product}
Let $b$ and $c$ be elements, and let $\pi$, $\pi'$, $\sigma$ and $\sigma'$ be automorphisms.
The pairs $(\pi(b),\sigma(c))$ and $(\pi'(b),\sigma'(c))$ are in the same orbit
if and only if the pairs $(\pi(\mathsf{supp}(b)),\sigma(\mathsf{supp}(c)))$ and $(\pi'(\mathsf{supp}(b)),\sigma'(\mathsf{supp}(c)))$ are in the same orbit.
\end{lemma}
\begin{proof}
We first show the forward implication.
Let $\rho$ be an automorphism that maps $(\pi(b),\sigma(c))$ to $(\pi'(b),\sigma'(c))$.
Since $\mathsf{supp}(\rho(\pi(b))) = \rho(\mathsf{supp}(\pi(b)))$,
the automorphism $\rho$ maps $\mathsf{supp}(\pi(b)) = \pi(\mathsf{supp}(b))$ onto $\mathsf{supp}(\pi'(b)) = \pi'(\mathsf{supp}(b))$.
Symmetrically, $\rho$ maps $\sigma(\mathsf{supp}(c))$ onto $\sigma'(\mathsf{supp}(c))$.

We now show the backward implication.
Let $\rho$ be an automorphism that maps $\pi(\mathsf{supp}(b))$ onto $\pi'(\mathsf{supp}(b))$ and $\sigma(\mathsf{supp}(c))$ onto $\sigma'(\mathsf{supp}(c))$.
The automorphism $\tau = \pi'^{-1}\rho\pi$ maps $\mathsf{supp}(b)$ onto itself.
Since $\tau$ preserves the order and $\mathsf{supp}(b)$ is finite,
$\tau$ fixes every element of $\mathsf{supp}(b)$.
Since $\mathsf{supp}(b)$ supports $b$, we get $\tau(b) = b$, that is, $\rho(\pi(b)) = \pi'(b)$.
Symmetrically, $\rho(\sigma(c)) = \sigma'(c)$.
So $\rho$ maps $(\pi(b),\sigma(c))$ to $(\pi'(b),\sigma'(c))$.
\end{proof}
%
\section{Polynomial rings}
%
The file \code{mon.rs} implements polynomials and basic operations on them.
A polynomial of type \code{Polynomial<F,X>} has coefficients of type \code{F},
which must be a field,
and variables of type \code{X}, which must be ordered.
For instance, \code{F} can be the field \code{F2} with two elements from \code{bits.rs},
or the field \code{BigRational} of rational numbers from the crate \code{num-rational}.
The standard way to define a polynomial is to first define the variables.
For instance, the code below show how to define a variable with type \code{String},
\[
\text{\code{let x : Polynomial<F2,String> = Polynomial::var("x".to\_string());}}
\]
\arka{this needs to clearer and more detailed}
Variables themselves are considered as polynomials.
Once we have defined variables $y$ and $z$ in a similar manner,
the polynomial $f = x + yz$ can be defined as
\[
\text{\code{let f = x.clone() + y.clone()*z.clone()}};
\] 
The \code{.clone()} is needed to ensure that the variables can be reused.
All the usual operations (addition, subtraction, multiplication, powering) of polynomials are defined in \code{mon.rs}.
Exponents are of type \code{Exponent}, which is \code{u32}.

The only important thing to note here is that monomials are compared lexicographically,
with the largest variable as the most significant one:
two monomials are compared by their exponents of the largest variable at which they differ.
For instance, $x^5 < y$ when $x < y$.

The file also implements the reduction of polynomials:
\code{reduce\_by} performs one reduction step of the leading term,
\code{remainder} repeats such steps until no reducer applies to the leading term,
and \code{normal\_form} computes the remainder of the division algorithm,
in which no monomial is divisible by the leading monomial of a reducer.
%
\subsection{Polynomial ideals}
%
Polynomials ideals are implemented in \code{ideal.rs}.
The function \code{buchberger} computes the reduced Gr\"{o}bner basis of the ideal generated by a list of polynomials.
It implements Buchberger's algorithm with Buchberger's two criteria as in \cite[Chapter 2, \S10]{CLO15},
with the normal selection strategy.

An ideal (\code{struct Ideal}) is stored by its reduced Gr\"{o}bner basis.
Since the reduced Gr\"{o}bner basis of an ideal is unique,
two ideals are equal if and only if their stored bases are equal.
The following operations are implemented:
membership (\code{contains}), inclusion (\code{is\_subset}), equality (\code{==}),
sum (\code{sum}) and product (\code{product}).

The next steps would be to implement quotients and intersections of ideals.
%
\section{Equivariant Buchberger's algorithm}\label{sec:eq buchberger}
%
The file \code{equivariant.rs} implements the algorithms of \cite[Sections 4 and 5]{GL25}
for polynomials whose variables are the atoms of an ordered structure $\struct$, such as $\dloQ$.
A list of orbits of polynomials stands for their union $H$,
and $I = \eqIdGen{H}$ is the equivariant ideal generated by $H$.
For a polynomial $f$, we write $\mathsf{supp}(f)$ for the set of its variables
and $\mathsf{LM}(f)$ for its leading monomial.
With the order of \code{mon.rs}, the largest variable of $f$ occurs in $\mathsf{LM}(f)$.

\paragraph{Reduction.}
A polynomial $f$ is only reduced by the elements $h$ of the orbits with $\mathsf{supp}(h) \subseteq \mathsf{supp}(f)$
(the \emphdef{support condition}).
A reduction step then does not introduce new atoms,
so reduction terminates, because the monomials over the finite set $\mathsf{supp}(f)$ are well ordered.
Following \cite{GL25}, an equivariant set $B$ is an \emphdef{equivariant Gr\"{o}bner basis} of $I$
if $\idGen{B} = I$ and every non-zero $f \in I$ has an element $h \in B$
with $\mathsf{supp}(h) \subseteq \mathsf{supp}(f)$ and $\mathsf{LM}(h)$ dividing $\mathsf{LM}(f)$.
Then $f \in I$ if and only if the normal form of $f$ is zero (\code{reduces\_to\_zero}):
reduction keeps $f$ in $I$, and a non-zero normal form in $I$ would be reducible.

\paragraph{Weak bases.}
The function \code{equivariant\_buchberger} adds the orbit of the normal form of every S-polynomial
that does not reduce to zero, until all S-polynomials reduce to zero.
Its result is a \emphdef{weak equivariant Gr\"{o}bner basis} \cite[Lemma 24]{GL25},
which is not always an equivariant Gr\"{o}bner basis.
For example, for $H$ the orbit of $x_a x_b - x_c$ with $a < b < c$ in $\dloQ$,
it stops with $5$ orbits, whose minimal elements are the orbits of
$x_c - x_a x_b$, $x_a^2 x_b - x_a x_b$ and $x_a x_b^2 - x_a x_b$.
However, $x_b - x_a \in I$ for $a < b$:
there are atoms $c < d < a$, and $x_a \equiv x_c x_d \equiv x_b$ modulo $I$.
No element $h$ of the $5$ orbits with $\mathsf{supp}(h) \subseteq \{a, b\}$
has a leading monomial dividing $x_b$, so $x_b - x_a$ does not reduce to zero.
Similarly, for $H$ the orbit of $x_a^2 - x_b$ with $a < b$,
the algorithm stops with the orbits of $x_b - x_a^2$, $x_a^4 - x_a^2$ and $x_a^6 - x_a^4$,
and the minimal elements are the orbits of $x_b - x_a^2$ and $x_a^4 - x_a^2$.
However, $x_a^2 - x_a \in I$:
for an atom $c < a$, we have $x_a \equiv x_c^2$ and $x_c^4 \equiv x_c^2$,
so $x_a^2 \equiv x_c^4 \equiv x_c^2 \equiv x_a$ modulo $I$.
The elements $h$ of the result with $\mathsf{supp}(h) \subseteq \{a\}$
are in the orbits of $x_a^4 - x_a^2$ and $x_a^6 - x_a^4$,
and neither $x_a^4$ nor $x_a^6$ divides $x_a^2$,
so $x_a^2 - x_a$ does not reduce to zero.
In both examples, the derivation needs atoms below the atoms of the polynomial,
which the support condition forbids.

\paragraph{The colouring.}
The function \code{equivariant\_groebner\_basis} uses the construction of \cite[Section 5]{GL25}.
Let $\mathcal{Y}$ be the set of atoms of $\struct$ with a colour, low or high,
where every low atom is smaller than every high atom (\code{struct Coloured<A>} in \code{coloured.rs}).
An automorphism $\pi$ of $\struct$ acts on $\mathcal{Y}$ by $\pi(c, a) = (c, \pi(a))$.
For a finite set $V$ of atoms, $\mathsf{col}_V$ maps the atoms in $V$ to low atoms and the others to high atoms,
and $\mathsf{forget}$ maps a coloured atom to its atom.
The function computes $\mathsf{forget}(B_\star)$,
where $B_\star$ is the weak basis computed by \code{equivariant\_buchberger}
from the coloured generators $\mathsf{col}_V(h)$, for $h \in H$ and $V \subseteq \mathsf{supp}(h)$.
By \cite[Lemma 28]{GL25}, this is an equivariant Gr\"{o}bner basis of $I$.
The idea of the proof is as follows.
Let $f \in I$ be non-zero and $V = \mathsf{supp}(f)$.
Then $\mathsf{col}_V(f)$ only has low variables,
and it has a decomposition by coloured generators whose low atoms all lie over $V$.
Since $B_\star$ is a weak basis, there is a decomposition with the same atoms
that contains a term $m h$, with $h \in B_\star$, such that $\mathsf{LM}(m h) = \mathsf{LM}(\mathsf{col}_V(f))$.
Every variable of $m h$ is at most the largest variable of $\mathsf{LM}(m h)$, which is low,
so every variable of $m h$ is low and lies over $V$.
Hence $\mathsf{forget}(h)$ satisfies $\mathsf{supp}(\mathsf{forget}(h)) \subseteq V$
and its leading monomial divides $\mathsf{LM}(f)$.

In the example above, \code{equivariant\_groebner\_basis} computes $10$ coloured orbits,
and the minimal elements of the result are the orbits of $x_b - x_a$ and $x_a^2 - x_a$.
So $I$ is the ideal of the polynomials that vanish at $(0, 0, \dots)$ and at $(1, 1, \dots)$.
The example \code{equivariant\_demo} runs both functions on several ideals.
The colouring multiplies the number of generators by up to $2^n$,
where $n$ is the size of their supports, so it is slower.
%
\paragraph{Finitely many variables.}
The file \code{trivial.rs} defines variables $x_1, x_2, \dots$ whose only automorphism is the identity.
For an ideal $I$ of $\poly{\F}{x_1, \dots, x_n}$,
an equivariant Gr\"{o}bner basis for the trivial group is a basis $B$ such that,
for every set $V$ of variables, the elements of $B$ in $\poly{\F}{X_V}$
form a Gr\"{o}bner basis of the elimination ideal $I \cap \poly{\F}{X_V}$.
Indeed, if $f \in I \cap \poly{\F}{X_V}$, the element $h$ given by the definition satisfies
$\mathsf{supp}(h) \subseteq \mathsf{supp}(f) \subseteq V$;
conversely, apply the property for $V = \mathsf{supp}(f)$.
By the Elimination Theorem \cite[Chapter 3, \S1]{CLO15},
an ordinary Gr\"{o}bner basis has this property for $V = \{x_1, \dots, x_k\}$,
because the largest variable of a monomial is its most significant one,
but not for every $V$:
for $I = \idGen{x_2 - x_1, x_1^2}$, the reduced Gr\"{o}bner basis is $\{x_2 - x_1, x_1^2\}$,
but $x_2^2 = (x_2 + x_1)(x_2 - x_1) + x_1^2 \in I$.
The function \code{elimination\_ideal} in \code{ideal.rs} computes the reduced Gr\"{o}bner basis of $I \cap \poly{\F}{X_V}$
with an elimination order, and the example \code{elimination\_comparison} compares three methods:
the ordinary Gr\"{o}bner basis,
the union of the reduced Gr\"{o}bner bases of the $2^n$ elimination ideals,
and the colouring.
%
\section{Dense meet-trees}\label{sec:trees}
%
The file \code{tree.rs} implements the atoms of the universal countable dense meet-tree with a lexicographic order,
written $\mathcal{T}_<$ in \cite[Example 12 and Remark 16]{GL25}.

\paragraph{Definitions.}
A \emphdef{meet-tree} is a partially ordered set $(A, \le)$ in which every two elements $a, b$
have a greatest lower bound $a \wedge b$, and the elements below any element form a chain.
A \emphdef{lexicographic order} on a meet-tree is a linear order $\preceq$ such that
(L1) $a \le b$ implies $a \preceq b$, and
(L2) for every $a$, the set $\setof{b}{a \le b}$ is an interval of $\preceq$.
We use the signature $(\wedge, \preceq)$, so the substructures are the subsets closed under $\wedge$.
The finite meet-trees with a lexicographic order form a Fra\"{i}ss\'{e} class,
and $\mathcal{T}_<$ is its Fra\"{i}ss\'{e} limit \cite[Remark 16]{GL25}.

\paragraph{The model.}
An atom is a sequence $(h_1, k_1, h_2, \dots, k_{n-1}, h_n)$ of rational numbers with $n \ge 1$ and every $k_i \neq 0$.
It is the point reached by going up the trunk to the height $h_1$, taking the branch with label $k_1$ there,
going up this branch to the height $h_2$, and so on.
Every branch is a copy of $(\Q, <)$, so it has no lowest point.
\begin{itemize}
\item $x \le y$ if $x$ without its last height is a prefix of $y$,
and the last height of $x$ is at most the height of $y$ at the same position.
\item For $u < x$, the \emphdef{direction} of $x$ at $u$ is $0$ if $x$ goes further up the branch of $u$,
and $k$ if $x$ takes the branch $k$ at $u$.
So the directions at every point are labelled by $\Q$.
\item $x \prec y$ if $x < y$, or if $x$ and $y$ are incomparable
and the direction of $x$ at $x \wedge y$ is smaller than that of $y$.
\end{itemize}

\begin{lemma}\label{lem:tree model}
The model, with $\wedge$ and $\preceq$, is isomorphic to $\mathcal{T}_<$. In particular, it is homogeneous.
\end{lemma}
\begin{proof}
The model is a countable meet-tree with a lexicographic order;
this follows from the definitions, and the tests check it on many points.
By Fra\"{i}ss\'{e}'s theorem \cite[Section 7.1]{Hod93},
it suffices to prove the extension property:
if $A \subseteq B$ are finite meet-trees with lexicographic orders,
then every embedding of $A$ into the model extends to $B$.
By induction on $|B \setminus A|$, we may assume that $B$ is generated by $A$ and one point $b \notin A$.
Let $u$ be the largest of the points $b \wedge a$ with $a \in A$, which form a chain below $b$.
Then $b \wedge a \le u \le b$, so $b \wedge a = u \wedge a$ for every $a \in A$, and $B = A \cup \{u, b\}$.
One of the following cases holds.
\begin{enumerate}
\item $u \in A$. Then $u < b$, and no point of $A$ lies in the direction of $b$ at $u$, by the choice of $u$.
The case is given by $u$ and by the place of the direction of $b$
among the directions at $u$ of the children of $u$ in $A$.
\item $u \notin A$. Two points of $A$ above $u$ have their meet in $A$, so it is not $u$:
the points of $A$ above $u$ lie in one direction at $u$, and they have a least element $q$.
So $u$ lies strictly between $q$ and its parent $p$ in $A$, or below $q$ if $q$ is the root of $A$,
and then $u \wedge a = q \wedge a$ for every $a \in A$ that is not above $q$.
Either $b = u$, or $b > u$ and the direction of $b$ at $u$ comes before or after that of $q$.
\end{enumerate}
In each case, the meets of $u$ and $b$ with the points of $A$ are determined,
and so is their order with the points of $A$, by (L1) and (L2).
The model has a point in each case, relative to the image of $A$:
every edge contains points, every point has points below it,
and at every point the directions are ordered like $\Q$,
so there are new directions before, between and after given ones.
The function \code{one\_point\_extensions} computes these points.
\end{proof}

If $A$ is closed under meets, the proof gives $|A|$ positions of a point $b \in A$,
$2|A| - 1$ positions in case 1 and $3|A|$ positions in case 2.
By homogeneity, these $6|A| - 1$ positions are the orbits of points
under the automorphisms that fix $A$.

\begin{lemma}\label{lem:tree orbits}
Let $S$ and $T$ be finite sets of atoms,
and let $f$ be the increasing bijection from $S$ onto $T$ for $\preceq$.
Some automorphism maps $S$ onto $T$ if and only if
$a \wedge b \le c \iff f(a) \wedge f(b) \le f(c)$ for all $a, b, c \in S$.
\end{lemma}
\begin{proof}
An automorphism that maps $S$ onto $T$ preserves $\preceq$, so it extends $f$,
and it preserves $\wedge$ and $\le$.
Conversely, assume the condition.
The meet of several points is the least of their pairwise meets,
because the points below any point form a chain.
So the closure of $S$ under meets is $\setof{a \wedge b}{a, b \in S}$.
Since $a \wedge b \le c \wedge d$ if and only if $a \wedge b \le c$ and $a \wedge b \le d$,
the map $g(a \wedge b) = f(a) \wedge f(b)$ is well defined and injective, and it preserves and reflects $\le$.
Hence it preserves $\wedge$.
It preserves $\preceq$: if $u < v$, then $g(u) < g(v)$ and we use (L1);
if $u$ and $v$ are incomparable, choose $a, b \in S$ with $u \le a$ and $v \le b$;
then $u \prec v$ if and only if $a \prec b$, by (L2),
and similarly $g(u) \prec g(v)$ if and only if $f(a) \prec f(b)$.
So $g$ is an isomorphism between finite substructures,
and it extends to an automorphism by \Cref{lem:tree model}.
\end{proof}

The function \code{in\_same\_orbit} checks the condition of \Cref{lem:tree orbits}.
The function \code{orbits\_in\_product(S, T)} fixes $S$ and adds the points of $T'$ one by one, in increasing order,
choosing each new point among the representatives of \code{one\_point\_extensions}
and keeping only the choices for which the points chosen so far have the type of the first points of $T$.
Each orbit of pairs is obtained exactly once,
because at each step the candidates represent the orbits under the automorphisms that fix the points already chosen.
The number of orbits of $k$-element sets is the number $c_k$ of plane trees with $k$ marked nodes
in which every unmarked node has at least two children:
$c_1, c_2, \dots = 1, 2, 8, 40, 224, \dots$, with generating function $F = x + 2F^2$.
The tests check these numbers,
and the number $\sum_k c_k \, k! / ((k-m)!\,(k-n)!\,(n+m-k)!)$ of orbits of pairs of an $n$-element set and an $m$-element set.

\paragraph{An example where the naive algorithm is not enough.}
Let $I$ be generated by the orbit of $x_b - x_a$ for $a < b$ in the tree order.
For incomparable $a$ and $b$, let $m = a \wedge b$.
Then $x_b - x_a = (x_b - x_m) - (x_a - x_m) \in I$.
The function \code{equivariant\_buchberger} returns the orbit of the generator alone,
and no element of this orbit has its support in $\{a, b\}$,
so $x_b - x_a$ does not reduce to zero.
The function \code{equivariant\_groebner\_basis} adds the orbit of $x_b - x_a$ for incomparable $a$ and $b$,
and $I$ is the ideal of the polynomials $f$ with $f(t, t, \dots) = 0$.
In $\dloQ$, this problem does not occur for this ideal, because any two atoms are comparable.
If $I$ is generated instead by $x_b - x_a$ for incomparable $a$ and $b$, the weak basis is enough:
for $a < b$, take $c$ incomparable to $a$ and $b$ and larger in the lexicographic order;
then the S-polynomial of $x_a - x_c$ and $x_b - x_c$ is $x_a - x_b$ up to sign.
For the ideal generated by $x_a^2 - x_b$ with $a < b$,
the weak basis misses both $x_a^2 - x_a$ and $x_b - x_a$ for incomparable $a$ and $b$,
and the ideal is generated by $x_b - x_a$ and $x_a^2 - x_a$, as in $\dloQ$.
Generators with three atoms, such as $x_a x_b - x_{a \wedge b}$ for incomparable $a$ and $b$,
are currently too slow: there are many orbits of pairs.
%
\begin{thebibliography}{9}
\bibitem{GL24}
Arka Ghosh and S\l{}awomir Lasota. Equivariant ideals of polynomials, LICS 2024.
%
\bibitem{GL25}
Arka Ghosh and Aliaume Lopez. Computability of equivariant
Gr\"{o}bner bases, 2025. arXiv:2507.08990.
%
\bibitem{Hod93}
Wilfrid Hodges. Model Theory. Cambridge University Press, 1993.
%
\bibitem{CLO15}
David A. Cox, John Little and Donal O'Shea. Ideals, Varieties, and Algorithms, 4th edition. Springer, 2015.
%
\end{thebibliography}
%
\end{document}
